% \chapterone


В современных информационных системах, будь то банковские хранилища или системы управления авиаперевозками, данные являются фундаментом для принятия любых решений. Однако между накопленными знаниями и конечным пользователем (аналитиком, менеджером или диспетчером) всегда стоит «барьер» в виде необходимости владения языком SQL. Традиционно любая выгрузка данных требует привлечения ИТ-специалиста, написания сложного запроса и его последующей оптимизации. В условиях динамично развивающейся отрасли, такой как авиация, задержка в получении информации может стать критичной.

Одной из главных проблем при работе с реляционными базами данных остается высокая сложность составления корректных и эффективных запросов. Например, демонстрационная база данных «Авиаперевозки» от компании Postgres Pro содержит множество связанных таблиц: от расписаний рейсов и бронирований до подробных схем компоновки бортов. Для получения простого ответа на вопрос «Какие самолеты чаще всего летали из Москвы в Санкт-Петербург в прошлом месяце с загрузкой более 80\%?» пользователю пришлось бы выполнить несколько сложных соединений (JOIN), учесть временные зоны и специфику хранения статусов рейсов. Вероятность ошибки в ручном SQL-запросе в таких условиях крайне высока.

Технология Text2SQL призвана радикально изменить этот процесс, превращая естественный язык в исполняемый программный код. Это не просто автоматизация написания скриптов, а создание «интеллектуального моста» между человеком и данными. Основные преимущества такого подхода заключаются в демократизации доступа к данным, снижением нагрузки на IT-специалистов, ускорением принятия решений и интерактивностью анализа, когда традиционный подход "запрос-ожидание-результат" заменяется на "вопрос-ответ-уточнение". В условиях, когда скорость и точность информации становятся конкурентными преимуществами, Text2SQL открывает новые горизонты для бизнеса, позволяя каждому сотруднику стать «информационным аналитиком» без необходимости изучать сложные языки запросов.

% ==============================================================
% Основная боль и эволюция подходов к решению проблемы NLIDB и Text-to-SQL
% ==============================================================

\section{Проблема NLIDB и эволюция подходов к Text-to-SQL}

Проблема создания интерфейсов взаимодействия с базами данных на естественном языке (Natural Language Interfaces to Databases, \texttt{NLIDB}) исследуется на протяжении нескольких десятилетий. Основная цель таких систем — преобразовать запрос пользователя в исполняемый \texttt{SQL}-код без потери семантического смысла. Эволюцию подходов к решению задачи \texttt{Text-to-SQL} можно разделить на несколько ключевых этапов, каждый из которых обладает своими преимуществами и ограничениями.

Исторически первыми системами трансляции текста в \texttt{SQL} были анализаторы, основанные на жестко заданных правилах и шаблонах. Системы извлекали ключевые слова из запроса пользователя и сопоставляли их с именами таблиц и столбцов. У этого подхода была крайне низкая масштабируемость. При малейшем изменении схемы базы данных или использовании синонимомов, система не могла корректно обработать запрос.

Далее для решения проблемы масштабируемости стали применяться методы семантического парсинга, где естественный язык сначала переводится в абстрактное логическое представление, а затем компилируется в \texttt{SQL}. Однако эти методы по-прежнему были ограничены в способности обрабатывать сложные запросы и требовали значительных усилий по ручной настройке для каждой новой предметной области. А так же требовался колоссальный объем ручной работы лингвистов и инженеров по созданию грамматик и словарей для каждой новой базы данных, что делало их непрактичными для широкого применения.

И уже с развитием глубокого обучения задача \texttt{Text-to-SQL} стала рассматриваться как задача машинного перевода, где входная последовательность — это текст, а выходная — \texttt{SQL}-код. Использовались рекуррентные нейронные сети (\texttt{RNN}) и архитектуры с механизмами внимания (\texttt{Attention}). Но была так же проблема: склонность к генерации синтаксически неверного кода и проблема «галлюцинаций» — модель могла выдумать несуществующие столбцы или таблицы, не ориентируясь на реальную схему базы данных.

А современный этап развития \texttt{NLIDB} характеризуется использованием предобученных языковых моделей на базе архитектуры \texttt{Transformer}. Эти модели обладают глубоким контекстным пониманием языка и структуры программного кода. Однако для промышленных систем, где критична точность (например, в авиационной аналитике), применение «чистых» генеративных моделей несет риски логических ошибок в сложных запросах с множественными соединениями (JOIN). Поэтому наиболее передовым подходом на сегодняшний день является использование архитектуры \texttt{RAG} (Retrieval-Augmented Generation) в связке с моделями понимания языка.

В рамках данного подхода, схема базы данных, метаданные и примеры сложных запросов векторизуются и помещаются в специализированную векторную базу данных. При поступлении пользовательского запроса система сначала извлекает наиболее релевантный контекст (структуру таблиц, описание специфичных полей). Обогащенный контекст передается в языковую модель для генерации итогового \texttt{SQL}-запроса, валидного для конкретной \texttt{СУБД}.



% ==============================================================
% Исследование существующих подходов к решению проблем NLIDB
% + сравнение аналогичных решений
% ==============================================================
\section{Сравнение существующих подходов и анализ аналогов}

\subsection{Сравнение существующих подходов}

\subsubsection{Современные направления в Text-to-SQL}

Современные подходы к решению задачи \textbf{Text-to-SQL} представляют собой переход от классических нейросетевых архитектур к сложным многоагентным системам и методам глубокого поиска, использующим возможности \textbf{больших языковых моделей (LLM)}.  \cite{hong2025survey} Текущее состояние области характеризуется разделением на две доминирующие парадигмы: \textbf{In-Context Learning (ICL)}, опирающуюся на продвинутый промпт-инжиниринг проприетарных моделей, и \textbf{Fine-tuning (FT)} открытых моделей для специфических диалектов и структур данных. \cite{li2024bird}

Одним из наиболее перспективных направлений является интеграция методов \textbf{поиска по дереву Монте-Карло (MCTS)}, реализованная в подходе \textbf{Alpha-SQL}. \cite{li2025alphasql} Этот метод формулирует генерацию SQL как задачу \textbf{итеративного поиска} в пространстве действий, где LLM выступает в роли модели действий, генерируя пошаговые рассуждения (Chain-of-Thought) для каждого узла дерева. Alpha-SQL использует самообучаемую функцию вознаграждения, основанную на \textbf{согласованности результатов выполнения} (self-consistency), что позволяет моделям среднего размера (например, 32B) превосходить такие гиганты, как GPT-4o, на сложных бенчмарках.

Для решения проблем \textbf{сложного математического анализа} и навигации по запутанным схемам предложен подход \textbf{SteinerSQL}. \cite{mao2025steinersql} В рамках этого фреймворка задача поиска оптимального пути соединения таблиц (join path) сводится к классической задаче \textbf{дерева Штейнера на графе схемы}. SteinerSQL объединяет математическую декомпозицию вопроса с детерминированными алгоритмами на графах, что гарантирует сохранение вычислительных зависимостей и минимизирует «шум» от избыточных соединений.

Проблема \textbf{надежности и галлюцинаций} в многоагентных системах эффективно решается с помощью технологии \textbf{consistency alignment} (выравнивание согласованности), представленной в \textbf{OpenSearch-SQL}. \cite{xie2025opensearch} Этот подход внедряет специализированный агент выравнивания, который гарантирует, что входные и выходные данные различных модулей (препроцессинг, экстракция, генерация) логически соответствуют друг другу, предотвращая накопление ошибок в цепочке агентов. Дополнительно в систему внедрен промежуточный язык \textbf{SQL-Like}, который упрощает структуру запроса перед генерацией финального синтаксиса.

В области \textbf{обучения с подкреплением} выделяется семейство моделей \textbf{Think2SQL}, использующее парадигму \textbf{RLVR} (Reinforcement Learning with Verifiable Rewards). \cite{papicchio2026think2sql} Вместо субъективных человеческих предпочтений модели обучаются на основе \textbf{проверяемых вознаграждений} от движка базы данных, используя плотные функции награды (Cell Precision, Cell Recall) для обеспечения стабильной конвергенции даже небольших моделей. Это позволяет модели «открывать» новые пути рассуждения, которые отсутствовали в исходных наборах данных для дообучения.

Для корпоративных систем, требующих \textbf{статистических гарантий корректности}, разработан фреймворк \textbf{RTS++}. \cite{chen2026rts} Он вводит понятие \textbf{execution entropy} (энтропия выполнения) для количественной оценки неопределенности модели и использует \textbf{конформное прогнозирование} (conformal prediction) для принятия решения о выполнении запроса или его передаче человеку на проверку. Такой подход критически важен в средах, где «тихие» ошибки (синтаксически верные, но семантически неверные запросы) недопустимы.

Наконец, стремление к \textbf{вычислительной эффективности} и защите данных привело к созданию \textbf{LitE-SQL} — легковесного фреймворка, демонстрирующего, что модели с 7B параметров могут достигать точности, сопоставимой с моделями в 175B–200B параметров. \cite{piao2025litesql} Успех этого подхода обеспечивается \textbf{векторным связыванием схемы} (Schema Linking) на основе контрастивного обучения и двухэтапным дообучением, что делает систему пригодной для локального развертывания.

\subsubsection{Современные направления в RAG для баз данных}

Современные направления в области \textbf{Retrieval-Augmented Generation (RAG)} для баз данных и их интеграция в задачи \textbf{Text-to-SQL} представляют собой переход от простых поисковых систем к сложным «Knowledge Runtime» средам. \cite{piao2025litesql} В 2026 году RAG уже не рассматривается как вспомогательный компонент, а становится ключевым слоем инженерии поиска, обеспечивающим актуальность знаний, соблюдение прав доступа и минимизацию стоимости эксплуатации за счет использования \textbf{малых языковых моделей (SLM)}. \cite{kmoseenk2026return}

Одним из основных направлений интеграции RAG в Text-to-SQL является \textbf{связывание схемы (Schema Linking)} на основе векторных баз данных. \cite{piao2025litesql} В отличие от ранних методов, современные системы, такие как \textbf{LitE-SQL}, хранят предобученные эмбеддинги метаданных схемы (названия таблиц, описания столбцов) в векторном хранилище и извлекают только релевантные элементы в ответ на запрос пользователя. \cite{piao2025litesql, piao2025arxiv} Это позволяет эффективно работать с огромными корпоративными схемами, содержащими тысячи столбцов, не переполняя контекстное окно модели и минимизируя шум от нерелевантных данных. \cite{ji2026benchmarks, piao2025litesql}

\textbf{Гибридный поиск (Hybrid Search)} и \textbf{реранкинг (Reranking)} признаны наиболее эффективными способами повышения точности извлечения. \cite{akzhankalimatov2026rag, kmoseenk2026return} Подход сочетает семантический векторный поиск, улавливающий общие смыслы, с лексическим поиском (BM25) для точного сопоставления технических терминов, идентификаторов и кодов ошибок. \cite{lin2025higress} Использование алгоритма \textbf{Reciprocal Rank Fusion (RRF)} для объединения результатов и последующее использование \textbf{кросс-энкодеров} для переранжирования топ-результатов позволяет повысить качество извлечения на 15–40\%. \cite{lin2025higress}

Развитие \textbf{агентного RAG (Agentic RAG)} привело к созданию модульных систем, где специализированные агенты (например, \textbf{Schema-RAG-Agent}) автономно планируют свои действия. \cite{chudinov2025agentic, xie2025opensearch} Агенты могут переформулировать запрос пользователя, решать, к какому источнику обратиться (документация, SQL, Slack), и проверять релевантность найденного контекста перед генерацией. \cite{akzhankalimatov2026rag, xie2025opensearch} Фреймворки типа \textbf{Corrective RAG (CRAG)} и \textbf{Self-RAG} добавляют уровень самокоррекции: если извлеченный контекст признан нерелевантным или амбивалентным, система инициирует повторный поиск или обращается к внешним инструментам (web search). \cite{akzhankalimatov2026rag, lin2025higress}

Инновационной парадигмой является \textbf{Table-Augmented Generation (TAG)}, которая объединяет Text-to-SQL и RAG для обработки запросов, требующих одновременно семантического анализа (мировых знаний) и точных вычислений в БД. \cite{biswal2024tag} Параллельно предложена концепция \textbf{Dual-Retrieval-Augmented-Generation (DRAG)}, формализующая рабочий процесс двойного извлечения схем таблиц и документов знаний перед генерацией SQL-запроса. \cite{ji2026benchmarks}

Для обеспечения надежности в корпоративных системах внедряются механизмы \textbf{выравнивания согласованности (Consistency Alignment)}, такие как в \textbf{OpenSearch-SQL}, и использование «энтропии выполнения» для оценки уверенности модели. \cite{xie2025opensearch, chen2026rts} Это превращает Text-to-SQL из вероятностной генерации в надежный сервис, способный воздержаться от ответа при высокой неопределенности. \cite{chen2026rts}


\subsection{Сравнение конкурентов}

В рамках разработки интеллектуального модуля был проведён сравнительный анализ программных решений, которые позволяют автоматизировать взаимодействие с реляционными базами данных на естественном языке. В качестве аналогов рассматривались наиболее востребованные инструменты в области \texttt{Text-to-SQL} и бизнес-аналитики — Vanna.AI, LangChain и Metabase.

\textbf{Vanna.AI} представляет собой специализированный Python-фреймворк с открытым исходным кодом, построенный на технологии \texttt{RAG}. Система обучается на метаданных, \texttt{DDL}-запросах и примерах пользовательских запросов, что обеспечивает высокую точность генерации для специфических схем баз данных. Однако, как показал анализ, отсутствует надёжный механизм самокоррекции: если сгенерированный \texttt{SQL}-запрос вызывает ошибку в \texttt{СУБД}, пользователю приходится править его вручную. К тому же возможности визуализации ограничены стандартными библиотеками и требуют дополнительного кода.

\textbf{LangChain} с его модулем \texttt{SQL Database Chain} — один из самых популярных фреймворков для создания приложений на базе больших языковых моделей. \cite{official_langchain} Он позволяет создавать цепочки и агентов, которые могут обращаться к схеме базы данных в реальном времени, и радует гибкостью настройки и обилием интеграций. Но у этого решения есть слабые места: стандартные цепочки LangChain для \texttt{SQL} начинают пробуксовывать на очень больших схемах (например, на демонстрационной базе «Авиаперевозки» с её множеством связей). Модель часто пытается передать в контекст всю схему целиком, что приводит к ошибкам или перерасходу токенов. Механизмы автоматического исправления ошибочных запросов в стандартных цепочках развиты слабо.

\textbf{Metabase} — это зрелая платформа для бизнес-аналитики, которая недавно обзавелась \texttt{AI}-помощником для генерации вопросов к данным. Она ориентирована на конечных бизнес-пользователей и создание дашбордов, предлагая отличную визуализацию «из коробки» и удобный интерфейс. Но Metabase — преимущественно облачное или серверное решение, что создаёт риски для конфиденциальности данных. Его \texttt{AI}-функционал слабо адаптирован для глубокой работы с \texttt{RAG}-системами на стороне клиента, а механизм самокоррекции практически отсутствует.

Спроектированный в рамках данной дипломной работы интеллектуальный модуль призван устранить перечисленные недостатки. Ключевое отличие — объединение трёх важных факторов в едином контуре. Во-первых, это автономность и приватность: использование локальных языковых моделей гарантирует, что данные авиапредприятия никогда не покинут внутренний периметр. Во-вторых, механизм самокоррекции: модуль реализует агентский подход, при котором в случае ошибки \texttt{СУБД} агент анализирует текст ошибки и автоматически пересобирает \texttt{SQL}-запрос до получения валидного результата. \cite{chaturvedi2025sql} В-третьих, глубокая проработка \texttt{RAG} для сложных схем. \cite{arxiv2410rag} В отличие от конкурентов, система оптимизирована для работы со специфическими связями базы данных Postgres Pro «Авиаперевозки» и использует векторное хранилище, чтобы извлекать только действительно нужные фрагменты схемы.

Сводный анализ функциональных возможностей исследуемых решений представлен в таблице~\ref{table:comparison}.

\begin{table}[h!]
    \centering
    \small   % уменьшаем шрифт (можно \footnotesize)
    \caption{Сравнение разрабатываемого решения с существующими инструментами}
    \label{table:comparison}
    % Устанавливаем относительную ширину столбцов:
    % первый столбец (функционал) шире, остальные — равномерно
    \begin{tabularx}{\textwidth}{|>{\hsize=1.2\hsize\raggedright\arraybackslash}X|
                                  |>{\hsize=0.95\hsize\centering\arraybackslash}X|
                                  |>{\hsize=0.95\hsize\centering\arraybackslash}X|
                                  |>{\hsize=0.95\hsize\centering\arraybackslash}X|
                                  |>{\hsize=0.95\hsize\centering\arraybackslash}X|}
    \hline
    \textbf{Функционал} & \textbf{Vanna.AI} & \textbf{LangChain SQL} & \textbf{Metabase} & \textbf{Мой продукт} \\
    \hline
    Генерация SQL (Text-to-SQL) & + & + & + & + \\
    \hline
    RAG для схемы БД & + & Ограничено & -- & + \\
    \hline
    Локальный запуск (Privacy) & + & + & -- & + \\
    \hline
    Механизм самокоррекции (Self-Heal) & -- & -- & -- & + \\
    \hline
    Авто-визуализация данных & $\pm$ & -- & + & Тесты \\
    \hline
    \end{tabularx}
\end{table} 




% ==============================================================
% Подробное обоснование целей и задач по дипломной работе 
% ==============================================================
\section{Обоснование цели и задач исследования}

На основе проведённого анализа существующих решений и выявленных пробелов в технологиях \texttt{Text-to-SQL} — в частности, отсутствия надёжных механизмов самокоррекции и ограничений при работе со сложными схемами данных — было сформулировано техническое обоснование цели и декомпозированных задач исследования.

\subsection{Обоснование цели работы}

Целью данной работы является разработка программного решения, которое устранит барьер между пользователями без технического образования и сложными реляционными базами данных. Создание полноценного \texttt{NLIDB}-интерфейса продиктовано необходимостью демократизации данных в крупных информационных системах, например, в сфере авиаперевозок. Это позволит конечным пользователям — менеджерам и аналитикам — получать оперативную информацию без ожидания выгрузок от ИТ-специалистов и тем самым снизит общую нагрузку на технический отдел предприятия.

Для достижения поставленной цели предстоит решить несколько взаимосвязанных задач.

Прежде всего, нужно проверсти оценку современных методов трансляции естественного языка в язык запросов. На сегодняшний день существует множество подходов и предобученных языковых моделей — как проприетарных, так и с открытым исходным кодом. Чтобы выбрать оптимальное «ядро» системы, потребуется провести сравнительный анализ их фичей.

Следующая задача — спроектировать \texttt{RAG}-архитектуру для работы с таблицами в схеме базы данных, с выявлением метаинформации и связей между таблицами. Передавать полную схему базы данных в контекстное окно языковой модели неэффективно и зачастую невозможно из-за ограничений на размер контекста. Архитектура \texttt{RAG} (Retrieval-Augmented Generation) позволяет динамически извлекать из векторного хранилища только те описания таблиц, полей и связей, которые действительно релевантны текущему вопросу пользователя. Такой подход гарантирует масштабируемость системы и сводит к минимуму риск галлюцинаций модели на сложных многотабличных базах данных.

Третья задача — разработка механизма самокоррекции ИИ-агента для повышения качества итогового ответа. Даже самые продвинутые языковые модели иногда генерируют \texttt{SQL}-запросы с синтаксическими или логическими ошибками. Подражая поведению человека-аналитика, ИИ-агент должен уметь перехватывать ошибки, возвращаемые СУБД, анализировать лог ошибки и автоматически переформулировать запрос. Внедрение такого цикла обратной связи (Self-Heal) заметно повысит автономность системы и долю успешно выполненных пользовательских запросов.

Четвёртая задача — провести оценку качества генерации после добавления в проект механизма самокоррекции. Любая дополнительная архитектурная надстройка требует валидации. Необходимо измерить метрики качества генерации, например точность выполнения запросов (Execution Accuracy), до и после внедрения механизма самокоррекции. Это позволит экспериментально подтвердить эффективность разработанного агентского подхода и обосновать его практическую ценность в рамках дипломного проекта.

Пятой задачей является сравнение существующих LLM в качестве \enquote{мозга системы}. Требуется протестировать отечественную разработку \texttt{Giga-chat 2} совместно с азиатскими моделями \texttt{Qwen3-Coder-Next} и \texttt{GLM-4.7}

Наконец, последней задачей является реализация микросервисной архитектуры. Интеллектуальный модуль должен обладать высокой отказоустойчивостью, гибкостью и независимостью компонентов. Разделение системы на микросервисы — сервис работы с \texttt{LLM}, сервис генерации эмбеддингов и векторного поиска, сервис-коннектор к реляционной базе данных — позволит независимо масштабировать наиболее нагруженные части системы, а также упростит интеграцию разработанного модуля в существующие корпоративные контуры предприятий.
